% Shared preamble of the HSEMHS mode documents (docs/physics/*.tex, docs/technical/*.tex).
\documentclass[10pt]{extarticle}
\usepackage[a4paper,margin=1.7cm]{geometry}
\usepackage[T1]{fontenc}\usepackage{lmodern}\usepackage{tabularx}
\usepackage{amsmath,amssymb,bm,enumitem,titlesec,xcolor,booktabs,array,tikz,graphicx,longtable}
\usepackage[hidelinks]{hyperref}
\titleformat{\section}{\large\bfseries\color{blue!50!black}}{\thesection.}{0.5em}{}
\titlespacing*{\section}{0pt}{9pt}{3pt}
\titleformat{\subsection}{\normalsize\bfseries}{\thesubsection}{0.5em}{}
\titlespacing*{\subsection}{0pt}{6pt}{2pt}
\setlist{nosep,leftmargin=1.3em}
\setlength{\parindent}{0pt}\setlength{\parskip}{4pt}
\renewcommand{\arraystretch}{1.15}
\newcommand{\B}{\mathbf B}\newcommand{\J}{\mathbf J}\newcommand{\F}{\mathbf F}\newcommand{\R}{\mathbf R}
\newcommand{\zh}{\hat{\mathbf z}}\newcommand{\Bo}{\mathbf B_{\rm obs}}\newcommand{\Lv}{\mathbf L}
\newcommand{\code}[1]{\texttt{#1}}
\newcolumntype{L}[1]{>{\raggedright\arraybackslash}p{#1}}
\newcolumntype{Y}{>{\raggedright\arraybackslash}X}
\newcommand{\doctitle}[2]{{\Large\bfseries #1}\\[2pt]{\small #2}\par\vspace{2pt}}
\pagestyle{plain}
\AtBeginDocument{\sloppy}

\begin{document}
\doctitle{mhs-poisson-ltau: implementation}{Technical notes on \code{run.mode = "mhs-poisson-ltau"}. HSEMHS, 26 Oct 2026 (after the embedding-API merge, 9d89191).}

\section{Running it}
{\footnotesize\begin{verbatim}
[mpiexec -n R] julia --project=. -t T scripts/run.jl config/modes/mhs-poisson-ltau.toml
\end{verbatim}}
Driver \code{run\_mhs\_poisson\_ltau} (\code{src/RunModes.jl}): read, \code{prepare\_embedded(COMM\_WORLD,
:mhs\_poisson\_ltau, \dots)}, \code{solve\_embedded!}, write, \code{free\_embedded!} (\code{src/Embedded.jl}); solver
\code{PoissonTau.solve\_tau}
(\code{src/PoissonTau.jl}); MPI layer \code{src/Slab.jl}; PCG and threaded vector updates from \code{PoissonZ}. Input and
\code{[tau\_grid]} keys as \code{column-hse-ltau}; in addition:

\begin{tabularx}{\linewidth}{L{4.8cm} L{3.4cm} Y}
\toprule
Key & Default & Meaning \\ \midrule
\code{poisson.logtau\_cut} & $-4$ & solved levels $\log\tau\ge$ cut (nearest level $k_c$) \\
\code{poisson.max\_steps} & 100 & Levenberg--Marquardt steps \\
\code{poisson.rtol}, \code{n\_consec} & $10^{-4}$, 3 & stop when $(\Phi_{\rm prev}-\Phi)/\Phi_{\rm prev}<$ rtol on \code{n\_consec} consecutive accepted steps (objective only) \\
\code{poisson.cg\_tol}, \code{cg\_maxit} & $10^{-3}$, 80 & PCG per trial \\
\code{poisson.dc\_cap\_m} & $3\times10^5$ & largest height change of a column per step (whole-step scaling) \\
\code{poisson.lm\_mu0} & 0 & initial damping \\
\bottomrule
\end{tabularx}

\section{Flow}
\code{column\_hse\_tau} (aligned, gauged; the current column integrator without a cap on the rise of $\ln P$ and with the
end-slope substep rule, see the \code{column-hse-ltau} notes) gives $u=\ln P$ and $z_0$ on all levels. If any column fails
there (pressure outside the table), the solve is not attempted and the driver raises an error. \code{solve\_tau} solves
for $u$ on levels $k>k_c$ and $c=z_{k_c}$ per column, with $u_{k\le k_c}$ and $z_k-z_{k_c}$ ($k<k_c$) fixed; a start state
outside the table returns \code{reason = :eos\_out\_of\_envelope} (driver error) instead of throwing. Then \code{gauge!}
(mean $z(\tau{=}1)=0$), \code{tau\_state!} ($\rho$, $n_e$, $\chi$ per node), \code{check\_fields!} (as in
\code{column-hse-ltau}), output in the \code{column-hse-ltau} layout.

Option checks in \code{prepare\_embedded}, identical on all ranks: \code{logtau\_cut} inside the axis; the cut level
$k_c\ge2$ (the height stencil at $k_c$ uses level $k_c-1$); at least 3 solved levels; \code{max\_steps}, \code{n\_consec},
\code{cg\_maxit} $\ge1$; \code{rtol} $>0$; $0<$ \code{cg\_tol} $<1$; \code{dc\_cap\_m} $>0$; \code{lm\_mu0} $\ge0$. The
distributed FFT plan (\code{SlabFFT}) is built once in \code{prepare\_embedded} and reused by every solve.

\section{State, heights and residuals (\code{evaluate!})}
Nodes $m=1\ldots n_1$ are levels $k=k_c\ldots n_\tau$. The unknown vector per column has $n_1$ slots: slot 1 $=c$, slots
$2\ldots n_1=u$ at $k>k_c$ (one array $(n_1,N_y,n_x)$, so FFTs and halos treat it uniformly). At each node the EOS gives
$\ln\rho$, $\partial\ln\rho/\partial\ln P$ and $\ln\chi$; $\partial\ln\chi/\partial\ln P$ by a central difference
($\varepsilon=10^{-4}$, the table returns no opacity derivative). With $a=\tau/\chi$, $q=\rho/P$:
\[
z_{k+1}=z_k-\tfrac12\Delta\ln\tau_k\,(a_k+a_{k+1}),\qquad z_{k_c}=c .
\]
Lorentz force at all levels (\code{lorentz\_tau!}): for each component $f$ of $\B$,
$\partial_zf=(f_{k-1}-f_{k+1})/(z_{k-1}-z_{k+1})$ (one-sided at the ends), $\partial_xf|_z=D_xf-D_xz\,\partial_zf$, likewise
$y$; $\J=\nabla\times\B/\mu_0$; $\Lv=\J\times\B$. Residuals with $W=gq-L_ze^{-u}$, $w=Wa$:
\[
R_x=-D_xu-W\,D_xz+L_xe^{-u},\qquad R_{z,f}=u_{f+1}-u_f-\tfrac12\Delta\ln\tau_f\,(w_f+w_{f+1}) .
\]
Weights from the start state (fixed): $\bar H_m$ global median of $1/(gq)$, $\overline{\Delta z}_f$ global median of the level
thickness, $w_c=\overline{\Delta z}_{\rm node}\bar H^2$, $w_f=\bar H_f\bar H_{f+1}/\overline{\Delta z}_f$;
$\Phi=\sum w_c(R_x^2+R_y^2)+\sum w_fR_z^2$.

\section{Jacobian}
Exact, including $\Lv$'s dependence on the heights (\code{Jop!}, adjoint \code{JTop!}):
\begin{itemize}
\item \textbf{Triangular operator.} $\delta z_m=\delta c+\sum_{f<m}\tfrac12\Delta\ln\tau_f\,(a\chi_u)_f\delta u_f+(a\chi_u)_{f+1}\delta u_{f+1}$ with
$\chi_u=\partial\ln\chi/\partial\ln P$: a cumulative sum down each column. Its adjoint is a reverse cumulative sum.
\item Vertical: $\delta R_{z,f}=\delta u_{f+1}-\delta u_f-\tfrac12\Delta\ln\tau_f(w'_f\delta u_f+w'_{f+1}\delta u_{f+1})$,
$w'=\partial w/\partial u=W'a-Wa\chi_u$, $W'=gq(\rho_u-1)+L_ze^{-u}$.
\item Horizontal: $\delta R_x=-D_x\delta u-W\,D_x\delta z+\ell_x\delta u$, $\ell_x=-D_xz\,W'-L_xe^{-u}$.
\item \textbf{Lorentz through the heights.} With $\zeta_m=(\delta z_{m-1}-\delta z_{m+1})/(z_{m-1}-z_{m+1})$ (the relative change of
the vertical spacing; $\delta z$ above the cut $=\delta c$): $\delta(\partial_zf)=-\partial_zf\,\zeta$,
$\delta(\partial_xf|_z)=-\partial_zf\,D_x\delta z+D_xz\,\partial_zf\,\zeta$. Hence
$\delta\J=\J_\zeta\zeta+\J_XD_x\delta z+\J_YD_y\delta z$ with
$\J_\zeta=(z_yB_{z,z}+B_{y,z},\,-B_{x,z}-z_xB_{z,z},\,z_xB_{y,z}-z_yB_{x,z})/\mu_0$, $\J_X=(0,B_{z,z},-B_{y,z})/\mu_0$,
$\J_Y=(-B_{z,z},0,B_{x,z})/\mu_0$, and $\delta\Lv=\delta\J\times\B$. It enters $\delta R_x$ through $e^{-u}(\delta L_x+z_x\delta L_z)$
and $\delta R_z$ through $e^{-u}a\,\delta L_z$: nine node-local coefficient arrays.
\item Checks (\code{MHS\_CHECK} in the script version): adjoint $\langle Jv,r\rangle=\langle v,J^\top r\rangle$ to $4\times10^{-16}$;
$Jv$ against central differences of $\R$ to $10^{-9}$ (with $\B$), $10^{-10}$ (without). With $\Lv$ held fixed the
directional error was 100\,\% in magnetic regions and the solve stalled after one step.
\end{itemize}

\section{Linear solve and preconditioner}
Normal equations $(J^\top WJ+E+\mu D)\,\delta=-J^\top W\R$ by PCG with global inner products; $E=10^{-2}\kappa^2D$,
$D=(\sum w_c\bar W^2,\ w_{c,2},\dots)$ (the damping shape; slot 1 for $c$, which no vertical equation sees).

\textbf{Preconditioner} (\code{setprec!}, \code{Minv!}): the operator with layer-median coefficients
($\overline{a\chi_u}$, $\overline{w'}$, $\bar W$) is, per horizontal mode, $M(s^2)=A+s^2B$ with $s^2$ as in
\code{mhs-poisson-z}; $A=G_z^\top W_fG_z+E+\mu D$ (vertical rows), $B=G_h^\top W_cG_h$, $G_h=-(I+\bar W\,S)$ the
horizontal symbol divided by $ik$ including the triangular matrix $S$. Both are dense $n_1\times n_1$ ($n_1=55$). Once per
step: Cholesky $A=LL^\top$, eigen-decomposition $L^{-1}BL^{-\top}=V\Lambda V^\top$, $Q=V^\top L^{-1}$; then
$M(s^2)^{-1}=Q^\top\,{\rm diag}(1/(1+s^2\lambda))\,Q$ for every mode: two complex GEMMs over all local modes of the
transposed FFT layout (\code{SlabFFT}, as in \code{mhs-poisson-z}), BLAS threaded during the solve.

\section{Steps}
Levenberg--Marquardt: for $\mu$ (start \code{lm\_mu0}) solve, scale the \emph{whole} step by
$\min(1,\ 2/\max|\delta u|,\ $\code{dc\_cap}$/\max|\delta c|)$, try up to 4 halvings on the true $\Phi$ (EOS status
checked); if none decreases $\Phi$, $\mu\leftarrow\max(4\mu,1)$ and re-solve (up to 8 times). After acceptance (source rule):
$\mu\leftarrow\mu/3$ if the accepted step length factor is exactly $\alpha=1$; unchanged if $\alpha\ge s_c/2$ ($s_c$ the
whole-step scale); $\times2$ otherwise; $\mu<10^{-3}\to0$. When the whole-step scaling is active ($s_c<1$), a full scaled
step ($\alpha=s_c$) therefore does \emph{not} lower $\mu$; whether that was intended is a code question (follow-up). Element-wise clamping
of $\delta c$ (the first version) turned the direction into an ascent direction where the step asked for $>1$\,Mm in a few
columns: $\Phi$ rose linearly with the step length.

\textbf{Termination.} (a) $(\Phi_{\rm prev}-\Phi)/\Phi_{\rm prev}<$ \code{rtol} on \code{n\_consec} consecutive accepted
steps (log line \code{converged: \dots}); (b) no damping level gives a decrease (\code{no decrease up to $\mu$ = \dots;
stopping}); (c) \code{max\_steps} reached (no message). The reason (\code{:converged}, \code{:no\_descent}, \code{:max\_steps}) is
returned in the status and written as the output attribute \code{stop\_reason}. Stopping uses the objective only. Two earlier rules are gone: a
single small step (it stopped one run after 17 steps while strong-field columns still moved 80\,km per step), and a
height shortcut (one small step plus $\max|\delta c|<100$\,m), removed on 24 Oct 2026 in line with the $\Phi$-only policy. The
logged $\max|\delta c|$ is the largest change of a column's \emph{cut-level} height in the step, not the change of its
$\tau=1$ height (which also depends on the changed layer thicknesses). The file driver still returns \code{converged}, exit 0, in all
three cases: read \code{stop\_reason} or the log (code follow-up: propagate the termination into the exit status). A met
rule does not certify small residuals.

\textbf{Stationarity diagnostic.} After the iteration, at the returned state and with $\mu=0$, one more system
$(J^\top WJ+E)\,\delta=-J^\top W\R$ is solved by PCG (diagnostic only; $u$, $z$ unchanged); logged as
\code{stationarity: Gauss--Newton decrement \dots} and written as \code{gn\_decrement} $=\delta\cdot(-J^\top W\R)/\Phi$, the
relative decrease a further undamped step would predict (0 at a stationary point); $\|J^\top W\R\|$ and the final $\mu$ are in
the embedded status. Option \code{rms\_atol} (embedded API only, default 0) accepts a start with $\sqrt{\Phi_0/W}\le$
\code{rms\_atol} in zero steps.

\section{Parallelisation}
$x$ slabs; halo exchange in $x$ for every $D_x$ (heights, $u$, $\B$, and in the Jacobian); global sums; exact global
medians for the weights and the preconditioner (bisection on order-preserving integer keys); the distributed FFT
(rfft in $y$, \code{Alltoallv} transpose, FFT in $x$) and GEMMs on the local $k_y$ slab. Threads in all node loops; the standalone driver raises BLAS to \code{min(32, nthreads)} around the solve and resets it to
1 (the embedded API leaves BLAS to the caller); FFTW plans are single-threaded, applied from Julia threads. The solve is
\emph{step-path sensitive}: on a $24\times18$ crop the same one-rank solve changes by O(1) in a few strong-field column heights
between 1 and 4 BLAS threads (different accepted LM steps), while $\Phi/\Phi_0$ agrees ($8.473\times10^{-3}$ vs
$8.474\times10^{-3}$).

The field values on the $\tau$ levels are used unchanged at the new heights; no discrete $\nabla\cdot\B=0$ is imposed (the
staggered construction of \code{full-mhs-z} needs a fixed geometry). Above the cut the levels stay at their column-HSE
pressures and relative heights and move rigidly with $c$; the global height offset is fixed only by the final gauge.

\section{Verification and performance}
\begin{itemize}
\item Automated tests: \code{test\_embedded.jl} (stored standalone outputs on two windows of the production input,
$24\times18$ and $11\times12$ columns, reproduced to $\le10^{-12}$; reason in \{converged, max\_steps, no\_descent\}; finite,
non-negative GN decrement; a balanced start accepted in zero steps; 1 vs 30 steps gives a larger decrement after 1);
\code{mpi/mpi\_embedded.jl}: on the converged $11\times12$ window (5 steps) 2, 3, 4 ranks agree with one rank to
$\le9\times10^{-11}$ (asserted $10^{-9}$); on the $24\times18$ window (24 steps) the differences are O(1) in a few heights
(step-path sensitivity, above) and only recorded.
\item $96^2$ magnetic crop (manual), ranks $\times$ threads $1\times8$, $1\times32$, $4\times8$,
consecutive-step rule: 37 steps each; $z(\tau{=}1)$ differences median 1\,m, 99th percentile 12--24\,m, max 150\,m
(weakly determined strong-field columns). Agreement to rounding only; identical trajectories across layouts are not
promised. After the removal of the height shortcut (commit 813a3af) the $1\times32$ run again took 37 steps.
\item Against the original script \code{scripts/poisson/poisson\_tau\_direct.jl} (single-step rule) on the same crop:
$z(\tau{=}1)$ differences median 3.6\,m, 99th percentile 95\,m, max 1\,km, because the two stopped one step apart while
some columns were still moving.
\item Full cube ($504\times504\times95$, 54 solved $\ln P$ levels + 1 height per column, $1.4\times10^7$ unknowns), with the
superseded single-step rule: one node, 96 threads, 17 steps, 14\,min, 11\,GB (commit b96bb4e); 8 ranks on 4 nodes, 8.3\,min,
3.1\,GB per rank, same summary statistics (commit 13666b7, log \code{out/modes/mhs-poisson-ltau.log} in the cluster
worktree \code{MHS\_modes}). These runs predate the removal of the $\ln P$ cap in the starting state; the current code has not been run on the full
cube. PCG reaches its
80-iteration cap in every step (relative residual $10^{-2}$--$2\times10^{-1}$): the next speed-up is a better preconditioner.
\end{itemize}

\section{Result on the production cube}
$z(\tau{=}1)$ relative to the cube mean, median [5\,\%, 95\,\%] in km, by $|\B|$ at $\tau=1$ (fixed-point = the superseded scheme;
24 Oct 2026, starting state still with the $\ln P$ cap, single-step stop rule):
\begin{center}\small
\begin{tabular}{lccc}
\toprule
$|\B|$ at $\tau=1$ & column HSE in $\tau$ & fixed-point scheme & \textbf{mhs-poisson-ltau} \\ \midrule
0--200\,G & $-38$ [$-85$, 187] & $-55$ [$-257$, 380] & $-10$ [$-35$, 73] \\
200--500\,G & $-68$ [$-88$, 58] & $+120$ [$-180$, 549] & $-26$ [$-52$, 24] \\
500--1000\,G & $-86$ [$-104$, $-36$] & $+183$ [$-198$, 587] & $-51$ [$-89$, $-17$] \\
1000--1500\,G & $-121$ [$-144$, $-92$] & $+314$ [$-117$, 563] & $-96$ [$-125$, $-63$] \\
1500--3000\,G & $-156$ [$-175$, $-140$] & $+402$ [58, 568] & $-130$ [$-167$, $-104$] \\
\bottomrule
\end{tabular}
\end{center}
Force imbalance after the solve, RMS in units of $\rho g$: 0.04--0.09 horizontal and 0.04--0.06 vertical for $\log\tau<0$
(start 0.3--0.8).

\section{Design history}
\begin{enumerate}
\item Fixed point (\code{scripts/poisson/poisson\_tau.jl}): map to a height grid, Poisson solve there, re-integrate the
$\tau$ columns, repeat. Each pass is a full Poisson solve (1\,h on the full cube), and it diverged where $\Lv$ carries the
gas: the heights were integrated with a pressure that was then replaced. A corrected variant (heights integrated with the Poisson
pressure) settled but lifted the strong-field columns instead of depressing them (median $+400$\,km above the cube mean for
1500--3000\,G).
\item Direct solve with $\Lv$ fixed in the Jacobian: exact without $\B$, stalled in magnetic regions.
\item Exact Jacobian + LM with whole-step scaling + the consecutive-step stop rule: the present mode.
\end{enumerate}
\end{document}
